\documentclass[a4paper,11pt]{ltjsarticle}
\usepackage[margin=25mm,headheight=16pt]{geometry}
\usepackage{amsmath,amssymb,amsthm,mathtools,bm,mathrsfs}
\usepackage{booktabs,longtable,array,enumitem}
\usepackage{xcolor}
\usepackage[unicode,colorlinks=true,linkcolor=blue!40!black,citecolor=blue!40!black,urlcolor=blue!45!black]{hyperref}
\usepackage{fancyhdr}
\usepackage{microtype}
\setlength{\parindent}{1em}
\setlength{\parskip}{0.25em}
\setlist{itemsep=0.2em,topsep=0.4em}
\allowdisplaybreaks[2]
\pagestyle{fancy}
\fancyhf{}
\fancyhead[L]{\small 3次超曲面のPoisson構造と非可換化}
\fancyhead[R]{\small 2026-11-09 / v6}
\fancyfoot[C]{\thepage}
\newcommand{\CC}{\mathbb C}
\newcommand{\ZZ}{\mathbb Z}
\newcommand{\PP}{\mathbb P}
\newcommand{\OO}{\mathcal O}
\newcommand{\cK}{\mathcal K}
\newcommand{\cP}{\mathcal P}
\newcommand{\cD}{\mathcal D}
\newcommand{\cC}{\mathcal C}
\newcommand{\cE}{\mathcal E}
\newcommand{\cF}{\mathcal F}
\newcommand{\cM}{\mathcal M}
\newcommand{\fm}{\mathfrak m}
\newcommand{\Perf}{\operatorname{Perf}}
\newcommand{\Ku}{\operatorname{Ku}}
\newcommand{\MF}{\operatorname{MF}}
\newcommand{\HH}{\operatorname{HH}}
\newcommand{\QGr}{\operatorname{QGr}}
\newcommand{\qgr}{\operatorname{qgr}}
\newcommand{\Dsg}{D_{\mathrm{sg}}}
\newcommand{\RHom}{\operatorname{RHom}}
\newcommand{\REnd}{\operatorname{REnd}}
\newcommand{\Hom}{\operatorname{Hom}}
\newcommand{\Ext}{\operatorname{Ext}}
\newcommand{\Ind}{\operatorname{Ind}}
\newcommand{\Def}{\operatorname{Def}}
\newcommand{\thick}{\operatorname{thick}}
\newcommand{\Spec}{\operatorname{Spec}}
\newcommand{\Spf}{\operatorname{Spf}}
\newcommand{\Id}{\operatorname{Id}}
\newcommand{\sym}{\operatorname{sym}}
\newcommand{\Cone}{\operatorname{Cone}}
\newcommand{\rank}{\operatorname{rank}}
\newcommand{\Pic}{\operatorname{Pic}}
\newcommand{\otL}{\mathbin{\otimes^{\mathbf L}}}
\newcommand{\DocDate}{2026年11月9日}
\newcommand{\DocVersion}{論文分割版 v6}
\theoremstyle{plain}
\newtheorem{theorem}{定理}[section]
\newtheorem{proposition}[theorem]{命題}
\newtheorem{lemma}[theorem]{補題}
\newtheorem{corollary}[theorem]{系}
\theoremstyle{definition}
\newtheorem{definition}[theorem]{定義}
\newtheorem{problem}[theorem]{問題}
\newtheorem{example}[theorem]{例}
\theoremstyle{remark}
\newtheorem{remark}[theorem]{注意}
\renewcommand{\proofname}{証明}
\hypersetup{pdfauthor={},pdfsubject={3次超曲面のPoisson構造・非可換量子化・rotation付き残余圏}}

\title{可換曲率付き対角の有限ホモトピー\\\large 全Artin底変更とkernel比較の残る条件}
\author{}
\date{2026年11月11日・研究草稿 v30}
\fancyhead[L]{\small 非可換三次超曲面の可換復元}
\fancyhead[R]{\small 2026-11-11 / v30}
\begin{document}
\maketitle
\begin{abstract}
非正則な可換Artin基底上でも、三次potentialのstabilized diagonalは、
有限自由matrix factorization上に恒等関手として作用する。
Koszulの差分商収縮と外積次数について有限な曲率付きホモトピー摂動により、
その評価射のホモトピー逆・ホモトピー・任意の底変更との両立を明示する。
これは相対kernel作用のunitに関する結果であり、
核Homから導来自然変換Homへの全忠実性、非可換rotation積の同定、
generic圏同値は未証明である。
さらに項別自由な周期複体がK-flatでない例により、
Artin底でbarの導来性を検証する必要を整理する。
\end{abstract}
% 2026-11-11 / v30
\section{可換曲率付き対角の有限ホモトピーとkernel作用}\label{dki:sec:identity}

本節は、論文Cの相対Artin節における条件(KF)のうち、
\emph{stabilized diagonalが有限rank graded matrix factorization上で恒等関手として作用すること}を、
任意の可換基底環上で有限のホモトピーにより証明する。
この結果は、kernel間のHom複体とdg双加群の導来自然変換複体が
擬同型であるという\emph{全忠実性}とは別である。
後者が必要な(BC)および非可換残余圏のrotation積は未完成として区別する。

以下の曲率付き摂動の原理はTu\cite[Appendix~A, Lemma~A.1]{DKITu}、
体上の積分変換・kernel比較はBallard--Favero--Katzarkov
\cite[Theorem~5.22]{DKIBFK}とDyckerhoff\cite[Section~5]{DKIDyckerhoff}
を参照する。特にBFKの定理は体上のdg圏の内部Homについての結果であり、
非正則Artin環へ仮定なしで移すものではない。
ここで直接示すのは\emph{objectwiseな恒等作用}とその底変更である。

\subsection{差分商による有限対角}

$R$を任意の可換環、$S_x=R[x_1,\ldots,x_n]$、
$S_y=R[y_1,\ldots,y_n]$、$S_e=R[x,y]$、
$W\in R[z_1,\ldots,z_n]$を斉次三次式とする。
$s_i=y_i-x_i$と置く。
$W_i=W(x_1,\ldots,x_i,y_{i+1},\ldots,y_n)$、$W_0=W(y)$、
$W_n=W(x)$と定め、
\begin{equation}
 a_i=\frac{W_i-W_{i-1}}{y_i-x_i}\in S_e
 \quad\Longrightarrow\quad
 \sum_{i=1}^n s_i a_i=W(x)-W(y).
 \tag{DKI1}\label{dki:eq:diff}
\end{equation}
右辺の割算は多項式の差分商であり、$s_i$を局所化しない。
$E=\bigoplus_iRe_i$、$K=S_e\otimes_R\bigwedge E$とし、
$\deg_{\rm coh}(e_i)=-1$、$\deg_{\rm int}(e_i)=1$、
$\deg t=(2,-3)$とする。
\begin{equation}
 d_0=\sum_i s_i\iota_{e_i^*},\qquad
 \delta=t\sum_i a_i(e_i\wedge-),\qquad
 d_K=d_0+\delta.
 \tag{DKI2}\label{dki:eq:koszul}
\end{equation}
すると$d_K^2=t(W(x)-W(y))$であり、$d_K$はcohomological degree一・
内部次数零を持つ。$K$をstabilized diagonalの明示的有限模型とする。

$M$を$S_y[t,t^{-1}]$上の有限自由graded MF、$d_M^2=tW(y)$とする。
$S_x$への射影に沿った、核の外積との畳込みの\emph{大きな}模型は
\begin{equation}
 \mathscr C(M)=K\otimes_{S_y}M,
 \qquad D_M=d_K\otimes1+(-1)^{|e|}\otimes d_M,
 \qquad D_M^2=tW(x).
 \tag{DKI3}\label{dki:eq:convolution}
\end{equation}
ここで$(-1)^{|e|}$は外積次数のparityであり、$S_x$-加群としては
$y$の多項式変数を残している。したがってこれを有限rankの
$S_x$-MFと\emph{そのまま}同一視しない。

\subsection{対角評価の明示的な強ホモトピー同値}

$M(x)$を係数置換$y_i\mapsto x_i$で得たMFとする。
$\pi_M:\mathscr C(M)\to M(x)$を、外積次数零への射影と
$y=x$での評価を合成して定める。

\begin{theorem}[有限対角のobjectwise恒等作用]\label{dki:thm:identity}
任意の可換基底環$R$、斉次三次式$W$、有限自由MF $M$に対し、
$\pi_M$はgraded MFの閉じた射であり、
\emph{明示的な有限和}で与えられる逆とホモトピーを持つ。
具体的に、多項式Koszul分解$(K,d_0)$の標準的な強変形収縮
\[
 (S_x\otimes_{S_y}M,0)\;
 \xleftarrow[\pi_0]{\ \ }(K\otimes M,d_0)
 \xrightarrow[\ ]{i_0}\;(S_x\otimes_{S_y}M,0)
\]
を$h_0$で表し、
\[
 p_M=\delta\otimes1+(-1)^{|e|}\otimes d_M(y),\qquad D_M=d_0+p_M
\]
とする（左端は$M$の自由基底を用いる同次加群としての表示）。
このとき
\begin{equation}\label{dki:eq:finite-hpl}
 \begin{split}
 i_M&=\sum_{j=0}^{n}(-h_0p_M)^j i_0,\\
 H_M&=\sum_{j=0}^{n}(-h_0p_M)^j h_0
 \end{split}
 \qquad\text{および}\qquad \pi_M=\pi_0
 \tag{DKI4}
\end{equation}
が定義され、次の恒等式を満たす：
\begin{equation}\label{dki:eq:identities}
 \begin{gathered}
 D_Mi_M=i_Md_M(x),\qquad
 \pi_MD_M=d_M(x)\pi_M,\qquad \pi_Mi_M=1,\\
 D_MH_M+H_MD_M=1-i_M\pi_M.
 \end{gathered}
 \tag{DKI5}
\end{equation}
以上の構成は任意の$R\to R'$への底変更と両立する。
$\mathbb G_m$の斉次次数と内部shiftを保ち、
三次斉次模型の$\mu_3$スカラー作用とも両立する。
\end{theorem}
\begin{proof}
一変数のKoszul分解$R[x,y]e_i\xrightarrow{s_i}R[x,y]$には
\[
 h_i(f)=e_i\frac{f(y_i)-f(x_i)}{y_i-x_i},\qquad
 h_i(e_if)=0
\]
という$R[x_i]$-線形な収縮がある。
その射影は$y_i=x_i$での評価、包含は$x_i$の多項式を
$y_i$に依らない式と見る写像である。
この収縮のtensor trickにより、$n$変数の$(K,d_0)$について
\[
 d_0h_0+h_0d_0=1-i_0\pi_0,\quad
 h_0^2=\pi_0h_0=h_0i_0=0
\]
を満たす収縮が得られる。例えば$n=2$では、外積次数零の
$f(y_1,y_2)$を
\[
 h_0(f)=
 e_1\frac{f(y_1,y_2)-f(x_1,y_2)}{y_1-x_1}
 +e_2\frac{f(x_1,y_2)-f(x_1,x_2)}{y_2-x_2}
\]
へ送り、$e_2f$を
$e_1\wedge e_2(f(y_1,y_2)-f(x_1,y_2))/(y_1-x_1)$へ送る。
この表示には$h$-進分母も基底環の非単元の反転もない。

$p_M$は外積次数を減少させない（$d_M$項は次数を保持し、
$\delta$項は一つ上げる）。$h_0$は外積次数を必ず一つ上げるので
\[
 (h_0p_M)^{n+1}=(p_Mh_0)^{n+1}=0.
\]
曲率は$D_M^2=tW(x)$であり、これはすべての演算子と可換する。
共通の中心曲率を持つ形式への有限homological perturbation
\cite[Appendix~A, Lemma~A.1]{DKITu}を適用すると、
式\eqref{dki:eq:finite-hpl}と\eqref{dki:eq:identities}が従う。
移送した微分は
\[
 \pi_0p_M(1+h_0p_M)^{-1}i_0=d_M(x)
\]
である。実際、増大射$\pi_0$は正の外積次数を消し、
$\delta$は外積次数を上げるため、唯一残るのは
$d_M(y)$を外積次数零に作用させ、$y=x$に評価した項である。
$\pi_0p_Mh_0=0$なので摂動後の射影は$\pi_0$に一致する。

すべての式は多項式の差分商と長さ高々$n+1$の有限和で書かれ、
$R\to R'$への底変更と可換する。内部次数は上の$e_i,t$の規約で
保存される。$\mu_3$のスカラー作用では、差分商、評価、外積シフトは
いずれも同次である。
\end{proof}

\begin{corollary}[自然な対角評価と有限Artinでの恒等性]
\label{dki:cor:unit}
$R_m=k[h]/(h^{m+1})$とすると、対角kernelの畳込みは
有限自由graded MFの各対象に、内部shiftと両立して
\emph{自然なホモトピー同値}を与える。
その特殊化は$R_0=k$上の同じ差分商対角評価である。
特に、(KF)に必要な\emph{対角の恒等作用とその底変更}は成立する。
\end{corollary}
\begin{proof}
次数零のMF射$f:M\to N$に対し、$1\otimes f(y)$は
畳込み側の閉射で、評価射について
\[
 \pi_N\circ(1\otimes f(y))=f(x)\circ\pi_M
\]
が厳密に成り立つ。高次のdg自然性は同じ式にKoszul符号を付けて得る。
各$\pi_M$は定理\ref{dki:thm:identity}により明示的ホモトピー逆を持つ。
特殊化は定理の底変更性による。
\end{proof}

\subsection{(KF)の残り：核作用の全忠実性}

\begin{remark}[恒等作用とkernel全忠実性の相違]\label{dki:rem:kf}
上の系は、相対stabilized diagonalが有限rank MF上で
\emph{恒等関手として作用すること}を示す。
これは\emph{相対kernel Homと導来自然変換Homの全忠実性}、すなわち
論文Cの相対Artin節で考えた写像
\[
 \operatorname{Hom}(\mathscr D_m,\mathscr D_m(r))
 \longrightarrow
 \RHom_{\mathcal T_m^e}(\Delta,\Delta(r))
 \tag{DKI6}\label{dki:eq:kf-map}
\]
が擬同型であることとは異なる。体上ではBFKのkernel定理が
対応する全忠実性を与えるが、Artin底では、同一enhancement上に
式\eqref{dki:eq:kf-map}の\emph{導来}模型を作り、二側barの値が
その模型に一致することを確認する必要がある。
定理\ref{dki:thm:identity}はそのためのunitの明示的な比較を提供するが、
(BC)や非可換rotation積の全同定は結論しない。
\end{remark}

\begin{example}[項別平坦性はK-flat性を意味しない]
\label{dki:ex:kflat}
$R=k[\varepsilon]/(\varepsilon^2)$とし、両方向に無限な複体
\[
 P:\quad\cdots\xrightarrow{\varepsilon}R
 \xrightarrow{\varepsilon}R\xrightarrow{\varepsilon}R\xrightarrow{\varepsilon}\cdots
\]
を考える。各項は自由で、$\ker(\varepsilon)=\varepsilon R$より
$P$は非輪状である。しかし$P\otimes_Rk$は零微分で
各次数のコホモロジーが$k$である。従って$P$はK-flatでない。
特に、(KF)の二側barの\emph{各項が平坦}であるという情報だけでは、
それが必要な導来tensorや導来双加群Homを計算することは保証されない。
一方、定理\ref{dki:thm:identity}は\emph{強ホモトピー同値}なので
この障害を回避し、任意の底変更で保存される。
\end{example}

\paragraph{再開すべき核心.}
指定生成子$K_m\oplus K_m(1)\oplus K_m(2)$の
適切な相対cofibrant双加群模型を作り、kernel作用の写像
\eqref{dki:eq:kf-map}を、体上BFKの全忠実性写像へ\emph{同じ鎖レベル}で
特殊化できることを確認する。単にgeneratorの特殊ファイバーでの
生成性や、対象毎の対角恒等作用だけから(BC)を導かない。

\paragraph{有限計算の範囲.}
\texttt{computations/verify\_diagonal\_identity.py}は、
一変数三次式と、混合二変数三次式
$W=y_1^3+y_2^3+2y_1^2y_2$を持つそれぞれrank2、rank4の
明示的MFについて、$t=1$で式\eqref{dki:eq:identities}、
$D_M^2=W(x)$、元のKoszul収縮を有理係数の多項式恒等式として検査する。
これは一般の$n$についての証明でも、(KF)の全忠実性の計算でもない。

\begin{thebibliography}{99}
\bibitem{DKITu}
J. Tu, \emph{Matrix factorizations via Koszul duality},
Compos. Math. \textbf{150} (2014), 1549--1578;
\href{https://arxiv.org/abs/1009.4151}{arXiv:1009.4151}, Appendix~A, Lemma~A.1.
\bibitem{DKIBFK}
M. Ballard, D. Favero and L. Katzarkov,
\emph{A category of kernels for equivariant factorizations and its implications for Hodge theory},
Publ. Math. Inst. Hautes \`Etudes Sci. \textbf{120} (2014), 1--111;
\href{https://doi.org/10.1007/s10240-013-0059-9}{doi:10.1007/s10240-013-0059-9},
Theorem~5.22, Remark~5.21.
\bibitem{DKIDyckerhoff}
T. Dyckerhoff, \emph{Compact generators in categories of matrix factorizations},
Duke Math. J. \textbf{159} (2011), 223--274;
\href{https://arxiv.org/abs/0904.4713}{arXiv:0904.4713}, Section~5.

\end{thebibliography}
\end{document}
